\documentclass[aspectratio=169]{beamer} % Formato 16:9 (más ancho)
\usetheme{Madrid}
\usecolortheme{default}
\usepackage[utf8]{inputenc}
\usepackage[spanish]{babel}
\usepackage{amsmath, amssymb}
\usepackage{hyperref}
\usepackage[spanish,ruled,linesnumbered]{algorithm2e}

% Reducir márgenes laterales (opcional)
\setbeamersize{text margin left=0.8cm, text margin right=0.8cm}

% Configurar espaciado en algoritmos (opcional)
\SetAlgoSkip{smallskip}

% Información actualizada
\title[Metastable Vacua]{Metastable Vacua}
\author{Brandon Marquez Salazar}
\institute{
    Maestría en Ingeniería Eléctrica (Instrumentación y Sistemas Digitales) \\
    4° Cuatrimestre \\
    \href{mailto:b.marquezsalazar@ugto.mx}{b.marquezsalazar@ugto.mx}
}
\date{}

\begin{document}

% Portada
\begin{frame}
    \titlepage
\end{frame}

% Índice
\begin{frame}{Contenido}
    \tableofcontents
\end{frame}

% Sección 1: Definición del problema
\section{Definición del problema / propósito del experimento}

\begin{frame}{Contexto: El programa Swampland}
    \begin{itemize}
        \item El programa Swampland busca criterios para separar teorías cuánticas de campos consistentes con gravedad cuántica (teoría de cuerdas) de aquellas que no lo son.
        \item Conjeturas relevantes:
        \begin{itemize}
            \item Inestabilidad de vacíos AdS no supersimétricos.
            \item Separación de escalas en AdS.
            \item Conjetura refinada de de Sitter (dS).
        \end{itemize}
        \item Se ha propuesto que estados no-BPS (clasificados por K-teoría) podrían evadir estas restricciones.
    \end{itemize}
\end{frame}

\begin{frame}{Propósito del experimento}
    \begin{block}{Objetivo principal}
        Construir vacíos dS (meta)estables mediante compactificaciones de teoría de cuerdas tipo IIB en variedades Kähler con torsión, identificando el conjunto mínimo de ingredientes en el potencial escalar efectivo.
    \end{block}
    \begin{block}{Objetivos específicos}
        \begin{itemize}
            \item Incorporar contribuciones del flujo \(F_5\) (por torsión) y de estados no-BPS \(\widehat{D5}\).
            \item Utilizar algoritmos de aprendizaje automático (ML) para explorar el espacio de parámetros.
            \item Analizar la relación con las conjeturas del Swampland.
        \end{itemize}
    \end{block}
\end{frame}

% Sección 2: Métodos empleados
\section{Definición de los métodos empleados}

\begin{frame}{Métodos: Algoritmo híbrido de Machine Learning}
    Combinación de dos algoritmos:
    \begin{enumerate}
        \item \textbf{Simulated Annealing (SA)}: búsqueda estocástica global.
        \item \textbf{Conjugate Gradient (CG)}: optimización determinística local.
    \end{enumerate}
    
    Función de error con restricciones físicas:
    \[
    \text{Error} = \sum_{i=1}^{N} \alpha_i \cdot \text{error}_i
    \]
    donde \(\alpha_i\) son pesos que penalizan regiones no deseadas del espacio de parámetros.
\end{frame}

\begin{frame}{Restricciones implementadas (errores)}
    \begin{itemize}
        \item \textbf{error\(_1\)}: \((\partial_j V)^2\) (búsqueda de puntos críticos)
        \item \textbf{error\(_2\)}: \(\|V\| - V\) (favorece \(V>0\), vacíos dS)
        \item \textbf{error\(_3\)}: \(\|\operatorname{tr} m_{ij}^2\| - \operatorname{tr} m_{ij}^2\) (traza positiva de la matriz de masas)
        \item \textbf{error\(_4\)}: \(\|\det m_{ij}^2 - \frac{1}{4}(\operatorname{tr} m_{ij}^2)^2\| + \det m_{ij}^2 - \frac{1}{4}(\operatorname{tr} m_{ij}^2)^2\) (autovalores positivos)
        \item \textbf{error\(_5\)}: \(\|A_{\widehat{D5}}\| - A_{\widehat{D5}}\) (contribución positiva de estados no-BPS)
    \end{itemize}
\end{frame}

\begin{frame}{Pseudocódigo: Simulated Annealing (SA)}
    \begin{algorithm}[H]
        \DontPrintSemicolon
        \caption{Búsqueda global con SA}
        Inicializar \(\phi_{\text{best}}\) aleatoriamente\;
        Fijar temperatura inicial \(T_0\) y esquema de enfriamiento\;
        \For{cada iteración \(i\)}{
            Generar nuevo punto \(\phi_{\text{new}}\)\;
            \(\Delta E = \text{Error}(\phi_{\text{new}}) - \text{Error}(\phi_{\text{best}})\)\;
            \If{\(\Delta E \le 0\)}{
                \(\phi_{\text{best}} = \phi_{\text{new}}\)\;
            }
            \If{\(\Delta E > 0\)}{
                Aceptar con probabilidad \(P = \exp\left(-\frac{\Delta E}{T_i}\right)\)\;
            }
            Actualizar \(T_{i+1}\)\;
        }
        \Return{\(\phi_{\text{best}}\)}
    \end{algorithm}
\end{frame}

\begin{frame}{Pseudocódigo: Conjugate Gradient (CG)}
    \begin{algorithm}[H]
        \DontPrintSemicolon
        \caption{Refinamiento local con CG}
        Inicializar \(\phi_0\) con la solución de SA\;
        \(g_0 = \nabla \text{Error}(\phi_0)\)\;
        \(d_0 = -g_0\)\;
        \For{\(k = 0,1,\ldots\) hasta convergencia}{
            Calcular \(\alpha_k\) por búsqueda lineal\;
            \(\phi_{k+1} = \phi_k + \alpha_k d_k\)\;
            \(g_{k+1} = \nabla \text{Error}(\phi_{k+1})\)\;
            \(\beta_k = \frac{g_{k+1}^T g_{k+1}}{g_k^T g_k}\)\;
            \(d_{k+1} = -g_{k+1} + \beta_k d_k\)\;
        }
        \Return{\(\phi_{\text{min}}\)}
    \end{algorithm}
\end{frame}

% Sección 3: Obtención de la función potencial
\section{Proceso para obtener la función potencial}

\begin{frame}{Reducción dimensional: Acción de tipo IIB}
    Acción en el marco de cuerda:
    \[
    S_{IIB} = S_G + S_\phi + S_{G_3} + S_{F_5} + S_{CS} + S_{\text{loc}}
    \]
    
    Compactificación en variedad Kähler \(\mathbb{X}_6\) con torsión:
    \begin{itemize}
        \item Permite contribución no trivial de \(F_5\) al potencial.
        \item Ciclos con torsión \(\Rightarrow\) estados no-BPS \(\widehat{D5}\).
    \end{itemize}
\end{frame}

\begin{frame}{Potencial escalar efectivo}
    Módulos: \(s = e^{-\phi}\) (axio-dilatón) y \(\tau = e^{-\phi} \mathcal{V}_6^{2/3}\) (Kähler)
    
    \[
    \mathcal{V}_{\text{eff}} = \frac{A_{H_3}s}{\tau^3} + \frac{A_{F_3}}{s\tau^3} + \frac{A_{F_5}}{\tau^4} + \frac{A_3\mathcal{N}_3}{\tau^3} + \frac{A_{\widehat{D5}}}{s^{1/2}\tau^{5/2}}
    \]
    
    donde:
    \begin{itemize}
        \item \(A_{H_3}, A_{F_3}\): flujos de tres-formas
        \item \(A_{F_5}\): contribución torsional de \(F_5\)
        \item \(A_3\mathcal{N}_3\): fuentes \(D3\) y \(O3\) (\(\mathcal{N}_3 = \mathcal{N}_{D3} - \frac{1}{2}\mathcal{N}_{O3}\))
        \item \(A_{\widehat{D5}}\): estados no-BPS
    \end{itemize}
\end{frame}

% Sección 4: Experimentos realizados
\section{Experimentos realizados}

\begin{frame}{Experimento 1: Búsqueda de vacíos AdS}
    \textbf{Configuración:}
    \begin{itemize}
        \item \(A_{\widehat{D5}} = 0\) (sin estados no-BPS)
        \item Condición: \(\mathcal{N}_3 < 0\) (más \(O3\) que \(D3\))
    \end{itemize}
    
    \textbf{Resultados:}
    \begin{itemize}
        \item 389 vacíos AdS estables encontrados
        \item No se encontraron vacíos dS
        \item Condición necesaria: flujos en más de un ciclo
    \end{itemize}
\end{frame}

\begin{frame}{Experimento 2: Búsqueda de vacíos dS}
    \textbf{Configuración:}
    \begin{itemize}
        \item Incluir \(A_{\widehat{D5}} > 0\) (estados no-BPS)
        \item Mantener \(\mathcal{N}_3 < 0\)
    \end{itemize}
    
    \textbf{Condiciones adicionales requeridas:}
    \begin{enumerate}
        \item Flujos \(H_3\) y \(F_3\) soportados en más de un ciclo
        \item Cota inferior para orientifoldos:
        \[
        \mathcal{N}_{O3} > 4\frac{\sqrt{A_{H_3}A_{F_3}}}{A_3} + 2\mathcal{N}_{\text{flux}}
        \]
    \end{enumerate}
    
    \textbf{Resultado:} Más de 170 vacíos dS estables
\end{frame}

% Sección 5: Resultados obtenidos
\section{Resultados obtenidos}

\begin{frame}{Ejemplos de vacíos dS encontrados}
    \begin{center}
    \begin{tabular}{c c c c c}
        \hline
        \(V_{\min}\) & \(\tau\) & \(s\) & \(A_{F_5}\) & \(A_{\widehat{D5}}\) \\
        \hline
        \(1.309\times10^{-6}\) & 0.00373 & 3.695 & 0.9769 & 0.4231 \\
        \(5.676\times10^{-6}\) & 0.00387 & 3.727 & 0.9755 & 0.4125 \\
        \(6.980\times10^{-5}\) & 0.00788 & 2.917 & 0.3022 & 0.1851 \\
        \hline
    \end{tabular}
    \end{center}
    (Valores parciales de la Tabla 1 del paper original)
\end{frame}

\begin{frame}{Mecanismo de uplifting AdS \(\rightarrow\) dS}
    Expansión para \(A_{\widehat{D5}}\) pequeño:
    \[
    \mathcal{V}_{\text{eff}} = \mathcal{V}_{AdS} + \frac{1}{s_0^{1/2}\tau_0^{5/2}} A_{\widehat{D5}} + \mathcal{O}(A_{\widehat{D5}}^2)
    \]
    
    \textbf{Consecuencias:}
    \begin{itemize}
        \item El mínimo se eleva de AdS a dS
        \item Los autovalores de la matriz de masas disminuyen
        \item El potencial se aplana \(\Rightarrow\) mayor probabilidad de inestabilidades
    \end{itemize}
\end{frame}

\begin{frame}{Relación con las conjeturas del Swampland}
    \begin{itemize}
        \item \textbf{AdS no supersimétricos:} metaestables por construcción (\(D_S W \neq 0\))
        \item \textbf{Separación de escalas:}
        \[
        m_{\text{mod}} R_{AdS} = \frac{3\sqrt{3}\,2\sqrt{A_{H_3}A_{F_3}} + A_3\mathcal{N}_3}{A_{F_5}} \gtrsim \mathcal{O}(1)
        \]
        \item \textbf{Conjetura dS:} vacíos obtenidos son muy planos, fácilmente desestabilizables
    \end{itemize}
\end{frame}

\begin{frame}{Conclusiones}
    \begin{itemize}
        \item La inclusión de \(F_5\) torsional y estados no-BPS \(\widehat{D5}\) permite construir vacíos dS clásicos estables en teoría de cuerdas tipo IIB.
        \item Se encontraron más de 170 vacíos dS mediante el algoritmo híbrido SA+CG.
        \item El uplifting produce potenciales muy planos, lo que sugiere inestabilidades potenciales.
        \item Los resultados son consistentes (en el límite) con las conjeturas del Swampland.
        \item Trabajo futuro: estudiar el rol de flujos torsionales de tres-formas y la cancelación de carga de K-teoría.
    \end{itemize}
\end{frame}

\end{document}
